\documentclass{article}
\usepackage{xcolor}
\usepackage{booktabs}
\usepackage{float}
\usepackage{graphicx}

\begin{document}
%\color{white}
%\pagecolor[rgb]{0.149,0.145,0.2}

\title{SVD Energy}
\date{}
\author{Denis Šmerda}
\maketitle

\section*{Introduction}
Hedging energy assets is an important aspect when trying to stabilize revenues.
Because a forward curve consists of dozens of individual contracts, it is incredibly inefficient to try and hedge each maturity date separately.
Instead, I need to identify the economic shocks that explain the volatility of these price changes.
In this project, I apply SVD to historical market data to extract these hidden risk factors, while taking a closer look at the numerical methods that can be applied.
Using this simplified market structure, I then implement least squares optimization to calculate the combination of exchange contracts needed to hedge an asset's risk.

\section*{Data}
The dataset I am working with provides hourly day ahead electricity prices in EUR/MWh for six western European countries.
The data is relatively clean, there are only a few observations missing in each country.
Moving forward, France is the country that will be analysed in this project, as the dataset is primarily focused on this country.
Also, there is strong seasonality in daily energy prices, but it differs depending on the weekday.
To make the models more accurate, I will focus only on business weekdays from now on, but the models could also be applied to only weekend data.

To address the missing values problem, simple linear interpolation across hours is used, rather than dropping the affected rows or applying more complex techniques such as splines.
After this change the dataset has 24 columns denoting all hours of a day, and 365 observations which stand for individual days ranging from 01/01/2022 to 31/12/2022.
Further exploring the data, I continue with summary statistics of individual hours over the year, presented in Table \ref{tab:Table 1}.
\begin{table}[H]
    \centering
    \caption{Summary statistics of individual hours}
      \begin{tabular}{lrrrr}
        \toprule
         & 3:00 & 9:00 & 15:00 & 21:00 \\
        \midrule
        mean & 233.8 & 360.3 & 291.2 & 346.0 \\
        std & 110.3 & 223.7 & 149.2 & 162.0 \\
        min & -0.0 & 22.0 & 28.6 & 15.0 \\
        25\% & 169.0 & 242.9 & 189.5 & 234.5 \\
        50\% & 203.8 & 304.3 & 237.2 & 290.6 \\
        75\% & 301.2 & 461.1 & 392.2 & 443.3 \\
        max & 585.0 & 2987.8 & 810.7 & 870.0 \\
        \bottomrule
      \end{tabular}
    \label{tab:Table 1}
\end{table}
Important discoveries of the summary statistics include the fact that prices reach negative values and there are some rather high observations.
The unusually high prices were recorded on the fourth of April from 8:00 to 10:00.
Various sources state that this was a supply shortage in a time of a huge demand, naturally increasing the price, thus it is not an observation error.
These extreme spikes are left unadjusted, as it can be interesting to see, how the SVD captures tail-risk events.

In Figure \ref{fig:raw_prices} I plot the midday prices over the whole year.
\begin{figure}[htbp]
    \centering
    \begin{minipage}{0.45\textwidth}
        \centering
        \includegraphics[width=\textwidth]{visual/raw_prices.png}
        \caption{Prices at Midday}
        \label{fig:raw_prices}
    \end{minipage}\hfill
    \begin{minipage}{0.45\textwidth}
        \centering
        \includegraphics[width=\textwidth]{visual/midday_differences.png}
        \caption{Midday price differences}
        \label{fig:midday_diff}
    \end{minipage}
\end{figure}
They are certainly very volatile and the volatility does not seem constant.
After running an ADF test on every column (hour) of this dataset, at the 5\% significance level there is not enough evidence to reject the test hypothesis, that there is no unit root, for 19 of them.
The following formula
$$
\Delta P_{day, hour} = P_{day, hour} - P_{day-1, hour}
$$
is applied to the whole dataset, resulting in Figure \ref{fig:midday_diff}.
It transforms the day-ahead prices to day-ahead price differences and is introduced to make all of the column time series stationary.
It has to transform the whole dataset, otherwise the SVD results would not make sense.
With this action the dataset is losing its first observation.
Logarithmic differences were not available in this dataset because of negative values.

Plot of differences in prices of each hour over the year can be seen in Figure \ref{fig:avg_hour}.
To deal with the seasonality issue, the following within transformation
$$
y_{day, hour}^{centered} = \Delta P_{day, hour} - \frac{1}{N} \sum_{day=1}^{N} {\Delta P}_{day, hour}
$$
is applied.
\begin{figure}[htbp]
    \centering
    \begin{minipage}{0.45\textwidth}
        \centering
        \includegraphics[width=\textwidth]{visual/avg_hour.png}
        \caption{Hour averages}
        \label{fig:avg_hour}
    \end{minipage}\hfill
    \begin{minipage}{0.45\textwidth}
        \centering
        \includegraphics[width=\textwidth]{visual/demeaned.png}
        \caption{Demeaned midday price differences}
        \label{fig:demeaned}
    \end{minipage}
\end{figure}
Midday absolute demeaned differences in day ahead prices throughout the year can be seen in Figure \ref{fig:demeaned}.

A portion of the summary statistics table after all preceding transformations is presented as Table \ref{tab: Table 2}.
\begin{table}[H]
  \centering
  \caption{Summary statistics of individual hours after transformation}
    \begin{tabular}{lrrrr}
      \toprule
      & 3:00 & 9:00 & 15:00 & 21:00 \\
      \midrule
      mean & -0.0 & 0.0 & 0.0 & 0.0 \\
      std & 50.4 & 232.6 & 61.6 & 54.4 \\
      min & -194.8 & -2538.8 & -337.4 & -191.0 \\
      25\% & -18.2 & -30.7 & -27.9 & -24.7 \\
      50\% & 1.1 & 0.7 & 2.2 & -0.4 \\
      75\% & 23.5 & 30.7 & 30.9 & 22.0 \\
      max & 248.1 & 2518.3 & 268.5 & 249.3 \\
      \bottomrule
    \end{tabular}
  \label{tab: Table 2}
\end{table}

\section*{SVD application}



\end{document}
